\documentclass[12pt,a4paper]{article}
\usepackage[utf8]{inputenc}
\usepackage[T1]{fontenc}
\usepackage{geometry}
\usepackage{amsmath, amssymb}
\usepackage{listings}
\usepackage{xcolor}
\usepackage{booktabs}
\usepackage{array}
\usepackage{hyperref}
\usepackage{graphicx}
\usepackage{tikz}
\usetikzlibrary{shapes.geometric, arrows.meta, positioning}

\geometry{left=2.5cm, right=2.5cm, top=2.5cm, bottom=2.5cm}

\definecolor{codebg}{rgb}{0.96,0.96,0.96}
\definecolor{codegreen}{rgb}{0.0,0.5,0.0}
\definecolor{codegray}{rgb}{0.5,0.5,0.5}
\definecolor{codepurple}{rgb}{0.58,0,0.82}
\definecolor{codeblue}{rgb}{0.0,0.0,0.8}

\lstdefinestyle{javastyle}{
    backgroundcolor=\color{codebg},
    commentstyle=\color{codegreen}\itshape,
    keywordstyle=\color{codeblue}\bfseries,
    stringstyle=\color{codepurple},
    numberstyle=\tiny\color{codegray},
    basicstyle=\ttfamily\small,
    breaklines=true,
    captionpos=b,
    keepspaces=true,
    numbers=left,
    numbersep=5pt,
    showspaces=false,
    showtabs=false,
    tabsize=4,
    language=Java,
    frame=single,
    framerule=0.5pt,
    rulecolor=\color{codegray}
}

\hypersetup{
    colorlinks=true,
    linkcolor=blue,
    urlcolor=cyan
}

\title{
    \Huge \textbf{Sorting Algorithms} \\[0.5em]
    \Large Implementation, Analysis \& Notes \\[0.5em]
    \large Data Structures \& Algorithms Reference
}
\author{Algorithm Study Guide}
\date{\today}

\begin{document}

\maketitle
\tableofcontents
\newpage

% ============================================================
\section{Introduction to Sorting}
% ============================================================

Sorting is the process of arranging elements in a particular order (ascending or descending). It is one of the most fundamental operations in computer science, as many algorithms rely on sorted data for efficient operation.

\subsection{Why Sorting Matters}
\begin{itemize}
    \item Binary search requires a sorted array.
    \item Database queries are faster on sorted data.
    \item Many algorithms (e.g., merge in merge sort) need sorted inputs.
    \item Sorting enables efficient data lookup, deduplication, and reporting.
\end{itemize}

\subsection{Complexity Overview}

\begin{table}[h!]
\centering
\renewcommand{\arraystretch}{1.4}
\begin{tabular}{|l|c|c|c|c|c|}
\hline
\textbf{Algorithm} & \textbf{Best Case} & \textbf{Average Case} & \textbf{Worst Case} & \textbf{Space} & \textbf{Stable?} \\
\hline
Quick Sort      & $O(n \log n)$ & $O(n \log n)$ & $O(n^2)$      & $O(\log n)$ & No \\
Merge Sort      & $O(n \log n)$ & $O(n \log n)$ & $O(n \log n)$ & $O(n)$      & Yes \\
Heap Sort       & $O(n \log n)$ & $O(n \log n)$ & $O(n \log n)$ & $O(1)$      & No \\
Insertion Sort  & $O(n)$        & $O(n^2)$      & $O(n^2)$      & $O(1)$      & Yes \\
\hline
\end{tabular}
\caption{Sorting Algorithm Complexity Summary}
\end{table}

\newpage

% ============================================================
\section{Quick Sort}
% ============================================================

\subsection{Concept}
Quick Sort is a \textbf{Divide and Conquer} algorithm. It selects a \textit{pivot} element and partitions the array into two sub-arrays: elements less than the pivot and elements greater than the pivot. These sub-arrays are then sorted recursively.

\subsection{Algorithm Steps}
\begin{enumerate}
    \item Choose a \textbf{pivot} element (e.g., last element — Lomuto scheme).
    \item \textbf{Partition}: rearrange the array so all elements $< \text{pivot}$ are on the left, $> \text{pivot}$ on the right.
    \item The pivot is now in its \textbf{correct final position}.
    \item Recursively apply Quick Sort to the left and right sub-arrays.
\end{enumerate}

\subsection{Partition (Lomuto Scheme)}
\[
i = \text{low} - 1
\]
For each $j$ from $\text{low}$ to $\text{high}-1$:
\[
\text{if } a[j] \le \text{pivot}: \quad i{+}{+}, \quad \text{swap}(a[i], a[j])
\]
Finally, swap pivot to position $i+1$.

\subsection{Recurrence Relation}
\[
T(n) = 2T\!\left(\frac{n}{2}\right) + O(n) \;\Rightarrow\; O(n \log n) \quad (\text{average case})
\]
\[
T(n) = T(n-1) + O(n) \;\Rightarrow\; O(n^2) \quad (\text{worst case: sorted input})
\]

\subsection{Complexity}
\begin{itemize}
    \item \textbf{Best / Average}: $O(n \log n)$
    \item \textbf{Worst}: $O(n^2)$ — occurs on already sorted or reverse-sorted arrays with naive pivot
    \item \textbf{Space}: $O(\log n)$ recursion stack
    \item \textbf{Stable}: No
\end{itemize}

\subsection{Java Implementation}
\begin{lstlisting}[style=javastyle, caption={Quick Sort — Java}]
public static void quickSort(int[] arr, int low, int high) {
    if (low < high) {
        int pi = partition(arr, low, high);
        quickSort(arr, low, pi - 1);
        quickSort(arr, pi + 1, high);
    }
}

private static int partition(int[] arr, int low, int high) {
    int pivot = arr[high];
    int i = low - 1;
    for (int j = low; j < high; j++) {
        if (arr[j] <= pivot) {
            i++;
            int temp = arr[i]; arr[i] = arr[j]; arr[j] = temp;
        }
    }
    int temp = arr[i + 1]; arr[i + 1] = arr[high]; arr[high] = temp;
    return i + 1;
}
\end{lstlisting}

\subsection{Key Notes}
\begin{itemize}
    \item Fastest in practice for large random datasets.
    \item Use \textbf{randomized pivot} to avoid worst-case $O(n^2)$.
    \item In-place sort (no extra array).
    \item Not stable — equal elements may be reordered.
\end{itemize}

\newpage

% ============================================================
\section{Merge Sort}
% ============================================================

\subsection{Concept}
Merge Sort is a stable, \textbf{Divide and Conquer} algorithm. It divides the array into halves, sorts each half recursively, then merges the sorted halves.

\subsection{Algorithm Steps}
\begin{enumerate}
    \item Find the midpoint: $\text{mid} = \lfloor (\text{left} + \text{right}) / 2 \rfloor$
    \item Recursively sort left half: $a[\text{left}..\text{mid}]$
    \item Recursively sort right half: $a[\text{mid}+1..\text{right}]$
    \item Merge both sorted halves into a single sorted array.
\end{enumerate}

\subsection{Recurrence Relation}
\[
T(n) = 2T\!\left(\frac{n}{2}\right) + O(n) \;\Rightarrow\; O(n \log n)
\]
By Master Theorem (Case 2): $T(n) = \Theta(n \log n)$ in \textbf{all cases}.

\subsection{Complexity}
\begin{itemize}
    \item \textbf{All cases (Best/Avg/Worst)}: $O(n \log n)$
    \item \textbf{Space}: $O(n)$ — auxiliary arrays for merge step
    \item \textbf{Stable}: \textbf{Yes}
\end{itemize}

\subsection{Java Implementation}
\begin{lstlisting}[style=javastyle, caption={Merge Sort — Java}]
public static void mergeSort(int[] arr, int left, int right) {
    if (left < right) {
        int mid = (left + right) / 2;
        mergeSort(arr, left, mid);
        mergeSort(arr, mid + 1, right);
        merge(arr, left, mid, right);
    }
}

private static void merge(int[] arr, int left, int mid, int right) {
    int[] L = Arrays.copyOfRange(arr, left, mid + 1);
    int[] R = Arrays.copyOfRange(arr, mid + 1, right + 1);
    int i = 0, j = 0, k = left;
    while (i < L.length && j < R.length)
        arr[k++] = (L[i] <= R[j]) ? L[i++] : R[j++];
    while (i < L.length) arr[k++] = L[i++];
    while (j < R.length) arr[k++] = R[j++];
}
\end{lstlisting}

\subsection{Key Notes}
\begin{itemize}
    \item Guaranteed $O(n \log n)$ regardless of input.
    \item Preferred when \textbf{stability is required}.
    \item Works well on \textbf{linked lists} (no random access needed during merge).
    \item Java's \texttt{Arrays.sort()} for objects uses a variant called \textbf{TimSort}.
    \item Extra $O(n)$ memory is a drawback for memory-constrained systems.
\end{itemize}

\newpage

% ============================================================
\section{Heap Sort}
% ============================================================

\subsection{Concept}
Heap Sort uses a \textbf{Binary Max-Heap} data structure. It first builds a max-heap from the array, then repeatedly extracts the maximum element and places it at the end.

\subsection{What is a Max-Heap?}
A complete binary tree where every parent node is \textbf{greater than or equal to} its children. Represented as an array:
\[
\text{parent}(i) = \lfloor (i-1)/2 \rfloor, \quad \text{left}(i) = 2i+1, \quad \text{right}(i) = 2i+2
\]

\subsection{Algorithm Steps}
\begin{enumerate}
    \item \textbf{Build Max-Heap}: Call \texttt{heapify} on all non-leaf nodes from bottom-up.
          This takes $O(n)$ time.
    \item \textbf{Extract-Max loop} (repeat $n-1$ times):
    \begin{enumerate}
        \item Swap root (max) with last element.
        \item Reduce heap size by 1.
        \item Call \texttt{heapify(root)} to restore max-heap property.
    \end{enumerate}
\end{enumerate}

\subsection{Heapify}
Heapify ensures the subtree rooted at index $i$ satisfies the max-heap property:
\[
\text{largest} = \text{argmax}(a[i],\; a[2i+1],\; a[2i+2])
\]
If $\text{largest} \ne i$, swap $a[i]$ with $a[\text{largest}]$ and recurse.

\subsection{Complexity}
\begin{itemize}
    \item \textbf{Build Heap}: $O(n)$
    \item \textbf{n Heapify calls}: $O(n \log n)$
    \item \textbf{Total (All cases)}: $O(n \log n)$
    \item \textbf{Space}: $O(1)$ — in-place!
    \item \textbf{Stable}: No
\end{itemize}

\subsection{Java Implementation}
\begin{lstlisting}[style=javastyle, caption={Heap Sort — Java}]
public static void heapSort(int[] arr) {
    int n = arr.length;
    // Build max-heap
    for (int i = n / 2 - 1; i >= 0; i--)
        heapify(arr, n, i);
    // Extract elements from heap
    for (int i = n - 1; i > 0; i--) {
        int temp = arr[0]; arr[0] = arr[i]; arr[i] = temp;
        heapify(arr, i, 0);
    }
}

private static void heapify(int[] arr, int n, int i) {
    int largest = i, left = 2*i+1, right = 2*i+2;
    if (left < n && arr[left] > arr[largest])   largest = left;
    if (right < n && arr[right] > arr[largest])  largest = right;
    if (largest != i) {
        int temp = arr[i]; arr[i] = arr[largest]; arr[largest] = temp;
        heapify(arr, n, largest);
    }
}
\end{lstlisting}

\subsection{Key Notes}
\begin{itemize}
    \item Guaranteed $O(n \log n)$ and $O(1)$ extra space — ideal when memory is tight.
    \item Not cache-friendly (poor locality of reference).
    \item Used in algorithms like \textbf{k-th largest element} (partial heap sort).
    \item Not stable — relative order of equal elements is not preserved.
\end{itemize}

\newpage

% ============================================================
\section{Insertion Sort}
% ============================================================

\subsection{Concept}
Insertion Sort builds the sorted array one element at a time. It picks each element and \textbf{inserts} it into its correct position among the already-sorted elements, shifting larger elements to the right.

Analogy: Similar to \textbf{sorting a hand of playing cards} — pick one card at a time and insert it in the right position.

\subsection{Algorithm Steps}
\begin{enumerate}
    \item Start from index $i = 1$ (index 0 is trivially sorted).
    \item Set $\text{key} = a[i]$.
    \item Compare \texttt{key} with elements at $j = i-1, i-2, \ldots$
    \item Shift elements $a[j] > \text{key}$ one position to the right.
    \item Insert \texttt{key} at the found position.
    \item Increment $i$ and repeat.
\end{enumerate}

\subsection{Complexity}
\begin{itemize}
    \item \textbf{Best Case}: $O(n)$ — already sorted; no shifts needed
    \item \textbf{Average / Worst Case}: $O(n^2)$
    \item \textbf{Space}: $O(1)$ — in-place
    \item \textbf{Stable}: \textbf{Yes}
\end{itemize}

\subsection{Java Implementation}
\begin{lstlisting}[style=javastyle, caption={Insertion Sort — Java}]
public static void insertionSort(int[] arr) {
    for (int i = 1; i < arr.length; i++) {
        int key = arr[i];
        int j = i - 1;
        while (j >= 0 && arr[j] > key) {
            arr[j + 1] = arr[j];
            j--;
        }
        arr[j + 1] = key;
    }
}
\end{lstlisting}

\subsection{Key Notes}
\begin{itemize}
    \item Excellent for \textbf{small} or \textbf{nearly-sorted} arrays.
    \item \textbf{Adaptive}: performs fewer comparisons if array is partially sorted.
    \item \textbf{Online}: can sort a list as elements arrive one by one.
    \item Used as the base case in \textbf{TimSort} (Java/Python's built-in sort) for small sub-arrays (size $< 64$).
    \item Stable: relative order of equal elements is preserved.
\end{itemize}

\newpage

% ============================================================
\section{Comparison \& When to Use Which}
% ============================================================

\begin{table}[h!]
\centering
\renewcommand{\arraystretch}{1.5}
\begin{tabular}{|>{\bfseries}l|p{10cm}|}
\hline
Quick Sort       & Use for large, random datasets. Fastest in practice. Use randomized pivot in production. \\
\hline
Merge Sort       & Use when stability is required, or for linked lists. Predictable $O(n \log n)$ always. \\
\hline
Heap Sort        & Use when $O(1)$ extra space is critical and stability isn't needed. \\
\hline
Insertion Sort   & Use for very small arrays ($n < 20$) or nearly-sorted data. Used as TimSort's base case. \\
\hline
\end{tabular}
\caption{Algorithm Selection Guide}
\end{table}

\section{Summary of Key Formulas}

\begin{align}
\text{Quick Sort (avg):}\quad & T(n) = O(n \log n) \\
\text{Merge Sort:}\quad        & T(n) = 2T(n/2) + O(n) = O(n \log n) \\
\text{Heap Sort:}\quad         & T(n) = O(n \log n),\; \text{Space} = O(1) \\
\text{Insertion Sort (best):}\; & T(n) = O(n) \quad (\text{sorted input})
\end{align}

\end{document}
